\section{True coordinate-ideal powers and genuine centre-order division memberships}

In the ACTUAL formal ring $S=\mathbb C[[X_1,\ldots,X_n]]$, fix ANY
actual subset $A$ of the original finite coordinate family. Put
$P_A=(X_i:i\in A)$, $w_i=1$ for $i\in A$ and $w_i=0$ otherwise,
and $w(\alpha)=\sum_iw_i\alpha_i$. This chapter gives the actual
coordinate-power meaning of the coefficient bounds in
\cite[Lemma~7.22 and Corollary~7.24, pp.~269--270]{bierstone-milman97-invent-math}.
The ideal, powers, formal products and scalar actions are the genuine
existing ones. No centre compatibility or strict-transform certificate
is manufactured by these formal algebra statements.

\begin{lemma}[Genuine weight-vanishing ideals contain the true coordinate powers]
\label{hir:actual-formal-coordinate-weight-vanishing-ideal}
For $q\ge0$, the actual set
$V_q=\{F\in S:[X^\alpha]F=0\text{ whenever }w(\alpha)<q\}$
is a genuine ideal of $S$. One has $P_A^q\subset V_q$, and for EVERY
actual exponent $\alpha$, the actual monomial $X^\alpha$ belongs to
$P_A^{w(\alpha)}$. In particular $w(\alpha)\ge q$ implies
$X^\alpha\in P_A^q$.
\lean{Hironaka.HironakaActualFormalCoordinatePowerDivision.actualFormalWeightVanishing,Hironaka.HironakaActualFormalCoordinatePowerDivision.actualFormalWeightVanishing_mul_mem,Hironaka.HironakaActualFormalCoordinatePowerDivision.actualCoordinateWeight,Hironaka.HironakaActualFormalCoordinatePowerDivision.actualFormalCoordinateIdeal,Hironaka.HironakaActualFormalCoordinatePowerDivision.actualFormalCoordinateIdeal_le_weightVanishing,Hironaka.HironakaActualFormalCoordinatePowerDivision.actualFormalCoordinateIdeal_pow_le_weightVanishing,Hironaka.HironakaActualFormalCoordinatePowerDivision.actual_monomial_mem_coordinate_weight_power}
\leanok
\end{lemma}
\begin{proof}
Zero and sums have the stated coefficient vanishing. For an ACTUAL
product $GF$ with $F\in V_q$, every genuine antidiagonal pair
$\beta+\gamma=\alpha$ satisfies $w(\gamma)\le w(\alpha)$.
Thus when $w(\alpha)<q$, the $F$ coefficient at $\gamma$ is zero,
and the entire genuine finite convolution vanishes. This proves ideal
closure for the actual ring multiplication.
Similarly, if $F\in V_p$ and $G\in V_q$, every pair with
$w(\alpha)<p+q$ has $w(\beta)<p$ or $w(\gamma)<q$, so actual
convolution gives $V_pV_q\subset V_{p+q}$. Each selected actual
variable belongs to $V_1$, so actual span and power induction give
$P_A^q\subset V_q$, including $q=0$.
Finally the genuine monomial is the actual finite product
$X^\alpha=\prod_iX_i^{\alpha_i}$. For selected $i$ the corresponding
factor lies in $P_A^{\alpha_i}$; for unselected $i$ it lies in the
actual whole ideal $P_A^0=S$. Multiplying these genuine ideal
memberships gives membership in $P_A^{\sum_{i\in A}\alpha_i}$,
exactly the claimed power. Monotonicity of genuine ideal powers gives
the final assertion. No infinite series division is used in this part.
\end{proof}

\begin{theorem}[The actual coordinate power has precisely the true coefficient kernel]
\label{hir:actual-formal-coordinate-power-coefficient-kernel}
\uses{hir:actual-formal-coordinate-weight-vanishing-ideal,hir:actual-formal-initial-exponent-coefficient-product,hir:actual-formal-vertex-division-identity,hir:actual-formal-vertex-division-supports}
For every genuine monomial order, the actual initial diagram of $P_A^q$
is exactly $\{\alpha:w(\alpha)\ge q\}$. For EVERY actual $F\in S$,
\[
 F\in P_A^q\quad\Longleftrightarrow\quad
 [X^\alpha]F=0\ \text{for every }w(\alpha)<q.
\]
These statements include empty and full coordinate subsets, $n=0$ and
$q=0$; the actual unselected coordinates are unrestricted.
\lean{Hironaka.HironakaActualFormalCoordinatePowerDivision.mem_actualInitialDiagram_coordinate_power_iff,Hironaka.HironakaActualFormalCoordinatePowerDivision.mem_actualFormalCoordinateIdeal_pow_iff,Hironaka.HironakaActualFormalCoordinatePowerDivision.actualFormalCoordinateIdeal_pow_eq_weightVanishing}
\leanok
\end{theorem}
\begin{proof}
An actual nonzero member of $P_A^q$ belongs to the preceding true
$V_q$. Its actual lowest nonzero coefficient therefore has weight at
least $q$. Conversely each exponent of weight at least $q$ has its
actual monomial in $P_A^q$ by the preceding lemma, and that monomial's
genuine lowest exponent is exactly the same exponent. This proves BOTH
actual diagram inclusions.
Power membership implies the displayed coefficient condition by
$P_A^q\subset V_q$. Conversely suppose $F\in V_q$ and apply
the already internally proved ACTUAL formal division for the WHOLE ideal
$P_A^q$. The finite sum of its genuine vertex products belongs to
$P_A^q\subset V_q$, so its actual remainder belongs to $V_q$ as
well. At every exponent of weight below $q$ that remainder's actual
coefficient therefore vanishes. At every exponent of weight at least
$q$, the just proved true diagram identity puts the exponent in the
actual diagram, so the genuine remainder support theorem again makes
its coefficient zero. Thus EVERY actual remainder coefficient is zero.
The WHOLE genuine division identity now places $F$ in $P_A^q$.
No coefficient kernel, quotient model or desired membership is a caller
premise, and no formal completion identification is needed.
\end{proof}

\begin{theorem}[Actual whole division quotients belong to the true centre-order ideal powers]
\label{hir:actual-formal-centre-power-whole-division-membership}
\uses{hir:actual-formal-coordinate-power-coefficient-kernel,hir:actual-formal-division-quotient-remainder-weight-bounds}
Let $I\subset S$ be ANY actual ideal and let its actual internally
constructed monic vertex series be $f_a$, for the fixed genuine monomial
order. Suppose each SAME actual $f_a$ belongs to $P_A^{w(a)}$.
For EVERY actual $F\in I\cap P_A^d$, its SAME genuinely constructed
division quotients satisfy
\[
 F=\sum_aQ_a(F)f_a,\qquad Q_a(F)\in P_A^{d-w(a)}.
\]
The natural-number subtraction convention handles $d<w(a)$. The stated
actual vertex power hypothesis is explicit; deriving it from genuine
admissible-centre geometry remains the separate transformation step.
\lean{Hironaka.HironakaActualFormalCoordinatePowerDivision.actual_formal_centre_power_whole_division}
\leanok
\end{theorem}
\begin{proof}
The BOTH-direction actual coordinate-power coefficient kernel turns
each genuine vertex power membership into exactly the actual coefficient
weight compatibility of the preceding division theorem. It similarly
turns the genuine numerator power membership into vanishing below
weight $d$. That theorem gives the actual WHOLE ideal member division
identity with its SAME quotients, and shows that a nonzero actual
coefficient of $Q_a(F)$ has weight at least $d-w(a)$. Equivalently
every actual coefficient below that weight is zero. The SAME true
coordinate-power kernel, now in the reverse direction, therefore
places the entire actual quotient series in $P_A^{d-w(a)}$, as claimed.
No quotient bound, division decomposition or ideal model is assumed.
\end{proof}
